\section{Introduction}
\label{sec:introduction}
Land-cover models are often trained on one mapping product and evaluated on another. Changes in imagery, sensor, resolution and geography can alter performance, while differences in the products' legends change what the learner is positively supervised to recognize \citep{tuia2016domain,bruzzone2009domain,yang2017harmonization}. If a target class is absent from the exact source labels, a near-zero target IoU is difficult to interpret without checking that absence. It is not sufficient evidence that visual domain shift alone caused the failure.

We call the finest software output category used in a mapping decision an \emph{operational leaf}; this is not a claim of a fully reconciled botanical taxonomy. Under a declared crosswalk, a target leaf may have exact source support, only coarse ancestral support, or no mapped support. These are supervision records, not mutually exclusive causes of transfer error. Exact-label absence also does not imply that a pretrained model has no representation of the concept, or that heterogeneous supervision cannot support it. It means that the specified source-training masks supply no positive exact-leaf examples, although the loss can still affect that leaf's predictions.

Label-space reconciliation and annotation conflicts are established concerns: MSeg reconciles taxonomies and pixel annotations, and UniSeg handles multi-dataset label shifts \citep{lambert2020mseg,kim2022uniseg}. Separately reporting common/private performance and examining negative transfer also have precedents \citep{choe2024bus,choe2025unida,ren2026sats}. Our question is narrower: what does a raw-mask support audit reveal about the interpretation of a specific transfer failure, and how much recovery and supported-leaf retention accompany the tested source mixture under fixed update and exposure controls?

LoveDA lacks a raw label assigned to the herbaceous output index in the frozen codebook, whereas FLAIR raw ID 10 supplies that training assignment. OpenEarthMap rangeland is a broader target-label proxy for this index, not a pure herbaceous reference (Supplementary Section S2). The primary CLIP-based experiment compares a LoveDA-only baseline with a LoveDA--FLAIR mixture while retaining the target replay population. Exposure and duplicate-crop controls constrain changes in LoveDA sampling and update representation. An archived native U-Net comparison checks the recovery/retention direction in another recorded model family on the same target. A recorded mismatched image--mask configuration supplies a separate, provenance-limited CLIP sensitivity. Auxiliary imagery and labels enter together, so none of these source-mixture comparisons isolates a label-only causal effect.

The contribution is an empirical audit and diagnostic design within this bounded setting. It connects a versioned operational crosswalk to raw-mask positive support; reports missing-leaf recovery beside source-supported retention under the specified controls; and examines recorded model-family and reversed-source extensions with their distinct dependence and provenance limits. The primary contrast is +0.3118 for rangeland IoU and -0.1011 for supported-leaf mIoU; archived U-Net records repeat the direction. The calibration adapter and U-Net are measurement instruments, not proposed algorithms. We make no first-discovery claim about label mismatch or negative transfer and claim no validated remedy for the observed trade-off.
